% Hodina 28, 1. ročník SPŠ. Rozsah: Rozpis_vyuky_MAT_1_rocnik.tex, řádek 28.

% Plán: 1_rocnik/tematicke_plany/SPS.pdf; společná projektová šablona.

% Učebnice: ucebnice/01_Zakladni_poznatky/Mocniny_a_odmocniny/Mocniny a Odmocniny.tex

% Původ úloh: P2 a procvičení a), b): učebnice, „Uprav výrazy“ a „Mocniny s nulovým exponentem“. Rezerva: přesný výraz z kostry učitele. Ostatní vlastní.

% Metodický zdroj: externi_zdroje/realisticky/Mocniny/05_Mocniny_s_celym_mocnitelem_I.pdf

% URL: http://www.realisticky.cz/ucebnice/01%20Matematika%20S%C5%A0/01%20Z%C3%A1kladn%C3%AD%20poznatky/06%20Mocniny/05%20Mocniny%20s%20cel%C3%BDm%20mocnitelem%20I.pdf

% Metodický zdroj: externi_zdroje/realisticky/Mocniny/06_Mocniny_s_celym_mocnitelem_II.pdf

% URL: http://www.realisticky.cz/ucebnice/01%20Matematika%20S%C5%A0/01%20Z%C3%A1kladn%C3%AD%20poznatky/06%20Mocniny/06%20Mocniny%20s%20cel%C3%BDm%20mocnitelem%20II.pdf

% Metodický zdroj: externi_zdroje/realisticky/Mocniny/07_Mocniny_s_promennou_v_mocniteli.pdf

% URL: http://www.realisticky.cz/ucebnice/01%20Matematika%20S%C5%A0/01%20Z%C3%A1kladn%C3%AD%20poznatky/06%20Mocniny/07%20Mocniny%20s%20prom%C4%9Bnnou%20v%20mocniteli.pdf

\documentclass[11pt,a4paper]{article}

% Spolecny styl pro ucitelske pripravy MAT (Work/MAT)
% Verze 1.1, 2026-10-23
\NeedsTeXFormat{LaTeX2e}
\ProvidesPackage{priprava}[2026/10/23 v1.1 Sablona ucitelskych priprav MAT]

% Zakladni typografie a cestina
\RequirePackage[T1]{fontenc}
\RequirePackage[utf8]{inputenc}
\RequirePackage[czech]{babel}
\RequirePackage{lmodern}
\RequirePackage{microtype}
\RequirePackage{amsmath,amssymb}
\RequirePackage{xcolor}
\RequirePackage{enumitem}
\RequirePackage[most]{tcolorbox}
\RequirePackage{tabularx,booktabs,array}
\RequirePackage{fancyhdr}
\RequirePackage{needspace}
\RequirePackage{geometry}
\RequirePackage[hidelinks]{hyperref}

\geometry{a4paper,top=17mm,bottom=18mm,left=23mm,right=23mm}
\definecolor{notegray}{gray}{0.39}
\setlength{\parindent}{0pt}
\setlength{\parskip}{2.5pt plus 1pt minus 0.5pt}
\setlist[itemize]{topsep=2pt,itemsep=1pt,parsep=0pt,leftmargin=1.45em}
\setlist[enumerate]{topsep=2pt,itemsep=2pt,parsep=0pt,leftmargin=1.65em}
\setcounter{secnumdepth}{0}
\renewcommand{\arraystretch}{1.12}

% Zahlavi: vlevo rocnik | globalni tema, vpravo cislo: nazev hodiny
% Zapati zustava zcela prazdne.
\pagestyle{fancy}
\fancyhf{}
\renewcommand{\headrulewidth}{0.35pt}
\renewcommand{\footrulewidth}{0pt}
\setlength{\headheight}{14pt}
\newcommand{\nastavPripravu}[4]{%
  \fancyhead[L]{\small #1\enspace\textbar\enspace #2}%
  \fancyhead[R]{\small #3: #4}%
}

% Minimalisticke cernobile boxy. Kazda definice zvlast.
\tcbset{commonlesson/.style={enhanced,breakable=false,colback=white,colframe=black,
  colbacktitle=white,coltitle=black,boxrule=0.45pt,
  arc=0.8mm,outer arc=0.8mm,left=3mm,right=3mm,top=1.6mm,bottom=1.6mm,
  toptitle=1mm,bottomtitle=0.7mm,fonttitle=\bfseries,
  attach boxed title to top left={xshift=2.5mm,yshift=-1.15mm},
  boxed title style={colback=white,colframe=white,boxrule=0pt,left=1mm,right=1mm}}}
\newtcolorbox{zadaniBox}[1][]{commonlesson,title={Zadání},#1}
\newtcolorbox{prikladBox}[1][]{commonlesson,title={Příklad},#1}
\newtcolorbox{definiceBox}[1][]{commonlesson,title={Definice},#1}

% Spolecny box pro zadani a jeho reseni, bez mezery mezi castmi.
% Pouziti:
%   \begin{prikladReseniBox}
%     Zadani ulohy.
%     \tcblower
%     \begin{align*} ... \end{align*}
%   \end{prikladReseniBox}
% Spodni cast ma jemny sedy podklad a delici linku; lze v ni sazet
% vice radku matematiky. Obe casti maji jen jeden spolecny ramecek.
\tcbset{solutionlesson/.style={
  commonlesson,
  enhanced jigsaw,
  breakable,
  colbacklower=black!3,
  segmentation style={solid,draw=black!25,line width=0.35pt},
  fontlower=\small\color{notegray},
  before lower={\noindent\textbf{Řešení:}\enspace},
  leftlower=3mm,rightlower=3mm
}}
\newtcolorbox{prikladReseniBox}[1][]{solutionlesson,title={Příklad},#1}
\newtcolorbox{zadaniReseniBox}[1][]{solutionlesson,title={Zadání},#1}

% Pomucky pro autora pripravy: sedive poznamky; \reseni pro starsi pripravy
\newcommand{\poznamka}[1]{\par\begingroup\small\color{notegray}\itshape
  \textnormal{\textbf{Poznámka pro učitele.}}\enspace #1\par\endgroup}
\newcommand{\reseni}[1]{\par\begingroup\small\color{notegray}
  \textbf{Očekávané řešení:}\enspace #1\par\endgroup}
\newcommand{\blok}[1]{\needspace{5\baselineskip}\subsection*{#1}}
\newcommand{\vyrok}[1]{\uv{#1}}


\nastavPripravu{1. ročník}{Mocniny}{28}{Nulový a záporný exponent}

\begin{document}

Z podílu $a^3:a^3=1$ při $a\ne0$ chceme zachovat pravidlo pro rozdíl exponentů. Podobně $a^2:a^3=1/a$. Rozšíření definice proto musí být slučitelné s již známým krácením.

\begin{definiceBox}[title={Definice},breakable=false]
Pro $a\ne0$ definujeme \textbf{nultou mocninu} předpisem $a^0=1$. Zápis $0^0$ zde nedefinujeme.
\end{definiceBox}


\begin{definiceBox}[title={Definice},breakable=false]
Pro $a\ne0$ a $n\in\mathbb N$ definujeme \textbf{mocninu se záporným celočíselným exponentem}: $a^{-n}=\frac1{a^n}$.
\end{definiceBox}


Pravidla pro mocniny nyní platí také pro celé exponenty, jsou-li všechny výrazy definovány. Záporný exponent znamená převrácenou hodnotu, nikoli záporný výsledek. Každý základ se záporným nebo nulovým exponentem musí být nenulový.

\begin{prikladReseniBox}[title={Příklad 1},breakable=false]
Vypočítej $2^{-3}$, $(-3)^{-2}$ a $\left(\frac23\right)^{-2}$.
\tcblower
$2^{-3}=\frac18$; $(-3)^{-2}=\frac19$; $\left(\frac23\right)^{-2}=\left(\frac32\right)^2=\frac94$.
\end{prikladReseniBox}


\begin{prikladReseniBox}[title={Příklad 2},breakable=false]
Uprav $(3a^7b^{-2}c^{-3})(4a^{-6}b^2c^{-1})$.
\tcblower
Původní zápis vyžaduje $a,b,c\ne0$. Potom $12a^{7-6}b^{-2+2}c^{-3-1}=\frac{12a}{c^4}$. Podmínku $b\ne0$ nesmíme po krácení vypustit.
\end{prikladReseniBox}


\blok{Samostatné procvičení}

\begin{zadaniReseniBox}[title={Procvičení},breakable=false]
\begin{enumerate}[label=\alph*)]
\item Uprav $2x^0$ a $(2x)^0$.
\item Uprav $\left(\frac{ab^{-3}}{c^4}\right)^{-2}$.
\item Uprav $\frac{x^{-2}y^3}{xy^{-1}}$.
\end{enumerate}
\tcblower
\begin{enumerate}[label=\alph*)]
\item Při $x\ne0$ dostáváme $2$ a $1$. Pro $x=0$ jsou oba původní zápisy nedefinované.
\item Pro $a,b,c\ne0$ vyjde $\frac{b^6c^8}{a^2}$. Vnější záporný exponent vyžaduje nenulovou celou závorku.
\item Pro $x,y\ne0$ je výsledek $x^{-3}y^4=\frac{y^4}{x^3}$.
\end{enumerate}
\end{zadaniReseniBox}


\poznamka{Před úpravou označ všechny původní nenulové základy a jmenovatele. Náročnou rezervu nech až po zvládnutí číselných příkladů a jednoduchých proměnných. Vyžaduj výsledek bez záporných exponentů tam, kde to zjednoduší čtení.}

\blok{Volitelná rezerva}

\begin{prikladReseniBox}[title={Rozšíření procvičení},breakable=false]
Pro $n\in\mathbb Z$ uprav $\frac{2^{3n+1}\cdot4^{2n-2}}{8^{n-1}}$.
\tcblower
Všechny základy jsou kladné, proto další omezení na celé $n$ nevzniká. Přepisem $4=2^2$, $8=2^3$ dostaneme $2^{3n+1+4n-4-3n+3}=2^{4n}$. Například pro $n=0$ je hodnota $1$.
\end{prikladReseniBox}


\end{document}
